% =============================================================================
% FS2D: Direct 1-Angle Rotation Correlation -- Mathematical Explanation
%
% Plain LaTeX content. Paste into your existing document setup.
% Uses the same macros/notation as the FS2D chapter:
%   \varmathbb{S}^{2}, \textbf{x}, \textbf{k}, align environments, \cite{...}
%
% This file explains (1) what the old full-SO(3) path computed,
% (2) what the new "direct" 1-angle path computes, (3) why they are
% mathematically equivalent for 2D, and (4) which old steps are no longer
% recomputed.
% =============================================================================

\subsection{Direct 1-Angle Rotation Correlation}
\label{sec:direct_1angle}

The rotation estimation described in Sec.\,\ref{sec:rotation_calculation} so far
forms the full three-dimensional SO(3) correlation of the two $\varmathbb{S}^{2}$
magnitudes, restricts it to the $\beta\approx 0$ slice, and finally collapses the
remaining $(\alpha,\gamma)$ plane into a one-dimensional rotation curve by
averaging over all $(\alpha,\gamma)$ pairs that share the same combined angle
$\theta=\alpha+\gamma$. We show in the following that, because a 2D rotation is a
single degree of freedom about the $z$-axis, this entire detour through the 3D
SO(3) correlation is unnecessary: the desired 1D rotation curve can be obtained
\emph{directly} from the spherical-harmonic coefficients of the two magnitudes,
without ever forming the SO(3) correlation, the Wigner-$d$ transform, or the
subsequent 2D-to-1D averaging. We call this the \emph{direct 1-angle} method.

\subsubsection{Shared front end}
\label{sec:direct_frontend}

Both the full-SO(3) method and the direct method share an identical front end up
to and including the two forward spherical transforms:
\begin{enumerate}
    \item Compute the 2D DFT $F(\textbf{k})$ of each scan $f(\textbf{x})$ and take
          the magnitude $|F(\textbf{k})|$.
    \item FFT-shift and normalize by the global maximum magnitude of both scans.
    \item Radially resample the shifted magnitude onto the
          $\varmathbb{S}^{2}$ grid $(\theta_j,\phi_k)$ as
          $f_{\varmathbb{S}^{2}}(\theta,\phi)$, see the resampling equations in
          Sec.\,\ref{sec:rotation_calculation}.
    \item Apply Contrast Limited Adaptive Histogram Equalization (CLAHE) and a
          1D Hamming window over $\phi$.
    \item Perform the forward spherical transform
          \texttt{FST\_semi\_memo} on each magnitude, yielding the
          spherical-harmonic coefficients
          $\hat{f}^{1}_{lm}$ and $\hat{f}^{2}_{lm}$ of the two scans.
\end{enumerate}
After step 5, both methods hold the same two coefficient arrays. They differ
only in how the rotation correlation is assembled from these coefficients.

\subsubsection{The full SO(3) path (reference)}
\label{sec:direct_so3path}

The full-SO(3) method combines the two $\varmathbb{S}^{2}$ coefficient sets into
SO(3) Fourier coefficients
\begin{align}
    \hat{C}^{l}_{mm'} \;=\; (-1)^{m-m'}\,\hat{f}^{1}_{l,-m}\,
                          \overline{\hat{f}^{2}_{l,-m'}},
    \label{eq:so3coef}
\end{align}
performs the inverse SO(3) Fourier transform to obtain the 3D correlation
$C(\alpha,\beta,\gamma)$ on the Euler-angle grid (ZYZ convention), and then
extracts the rotation curve as follows. Because the projected magnitudes are
real, the only slice that carries the 2D rotation is the one with $\beta\approx
0$, where the inverse transform is sampled at the first Gauss--Legendre node
$\beta_{1}=\pi/(2N)$. On this slice, every pair $(\alpha,\gamma)$ with the same
sum $\theta=\alpha+\gamma$ corresponds to the same physical 2D rotation, so all
such pairs are averaged to produce the 1D curve
\begin{align}
    C_{\text{1D}}(\theta) \;=\; \text{avg}_{\alpha+\gamma=\theta}\;
        C(\alpha,\beta_{1},\gamma),
    \label{eq:so3_collapse}
\end{align}
which is then normalized to $[0,1]$ and searched for peaks. The dominant cost of
this path is the inverse SO(3) transform, which scales as
$O(N^{3})$ in the bandwidth $B=N/2$.

\subsubsection{The direct 1-angle path}
\label{sec:direct_directpath}

The direct method skips the SO(3) combine and the inverse SO(3) transform
entirely. From the two spherical-harmonic coefficient sets it first computes, for
every order $m$, the per-order inner product
\begin{align}
    P_{m} \;=\; \sum_{l\geq |m|} \hat{f}^{1}_{lm}\,
               \overline{\hat{f}^{2}_{lm}},
    \label{eq:Pm}
\end{align}
which is a single sum over the degree $l$ and costs $O(B^{2})$ in total over all
$m$. The 1D rotation correlation is then evaluated directly as a Fourier series
in the combined angle $\theta$:
\begin{align}
    C(\theta) \;=\; \sum_{m} P_{m}\,e^{-im\theta},
    \qquad \theta_{k}=\frac{2\pi k}{N},\;\; k=0,\dots,N-1.
    \label{eq:direct_corr}
\end{align}
The cost of evaluating Eq.\,\ref{eq:direct_corr} on all $N$ angles is $O(NB)$,
after which the curve is normalized to $[0,1]$ and searched for peaks with the
same peak detection as before. No Wigner-$d$ functions, no SO(3) coefficients,
and no 2D-to-1D averaging are formed.

\subsubsection{Why the direct method is mathematically equivalent}
\label{sec:direct_equivalence}

The equivalence follows from the structure of the SO(3) correlation
\begin{align}
    C(g) \;=\; \sum_{l}\sum_{|m|\leq l}\sum_{|m'|\leq l}
              \hat{f}^{1}_{lm}\,\overline{\hat{f}^{2}_{lm'}}\,
              (-1)^{m-m'}\,D^{l}_{mm'}(g),
    \label{eq:so3_corr_full}
\end{align}
with the Wigner-$D$ function in ZYZ Euler angles
\begin{align}
    D^{l}_{mm'}(g(\alpha,\beta,\gamma))
        \;=\; e^{-im\alpha}\,d^{l}_{mm'}(\beta)\,e^{-im'\gamma},
\end{align}
where $d^{l}_{mm'}(\beta)$ is the Wigner small-$d$ matrix. A 2D rotation about
the $z$-axis is a single degree of freedom, so only the slice $\beta=0$ of the
SO(3) correlation is relevant. At $\beta=0$ the small-$d$ matrices collapse to
the identity,
\begin{align}
    d^{l}_{mm'}(0) \;=\; \delta_{mm'},
\end{align}
so the $(m,m')$ double sum in Eq.\,\ref{eq:so3_corr_full} collapses to a single
sum over $m$:
\begin{align}
    C(\alpha,0,\gamma)
        &\;=\; \sum_{l}\sum_{m}
              \hat{f}^{1}_{lm}\,\overline{\hat{f}^{2}_{lm}}\,
              e^{-im\alpha}\,e^{-im\gamma} \\[2pt]
        &\;=\; \sum_{m}
              \underbrace{\Bigl(\sum_{l}\hat{f}^{1}_{lm}
              \overline{\hat{f}^{2}_{lm}}\Bigr)}_{P_{m}}\,
              e^{-im(\alpha+\gamma)} \\[2pt]
        &\;=\; \sum_{m} P_{m}\,e^{-im\theta},
        \qquad \theta=\alpha+\gamma.
    \label{eq:collapse}
\end{align}
Equation\,\ref{eq:collapse} is exactly Eq.\,\ref{eq:direct_corr}. The direct
method is therefore not an approximation of the full SO(3) method but its
analytic $\beta=0$ reduction: it computes precisely the quantity the full method
recovers only after assembling the 3D correlation, taking a $\beta$-slice, and
averaging over $(\alpha,\gamma)$.

In fact, the direct method is \emph{more} exact than the full-SO(3) reference in
one respect: the full method evaluates the $\beta$ slice at the first
Gauss--Legendre node $\beta_{1}=\pi/(2N)$, where
$d^{l}_{mm'}(\beta_{1})\neq\delta_{mm'}$ exactly, and then averages over
$(\alpha,\gamma)$; the direct method evaluates $\beta=0$ analytically, where the
collapse $d^{l}_{mm'}(0)=\delta_{mm'}$ holds exactly. In practice the two
normalized 1D curves agree to within numerical precision: a side-by-side
comparison on real sonar scans shows a maximum pointwise difference below $1\%$
after $[0,1]$ normalization.

\paragraph{Sign convention and the evenness of the rotation curve.}
The full-SO(3) path labels the rotation axis with $\theta_{\text{SO3}}=
-(\alpha+\gamma)$, whereas the direct path labels it with $\theta_{\text{direct}}
=+\alpha=+\theta$. These differ by a sign. Because the projected magnitudes are
real, the spherical-harmonic coefficients satisfy
$\hat{f}_{l,-m}=(-1)^{m}\overline{\hat{f}_{lm}}$, which implies
$P_{-m}=\overline{P_{m}}$ and hence
\begin{align}
    C(-\theta) \;=\; C(+\theta).
    \label{eq:evenness}
\end{align}
The 1D rotation correlation is an \emph{even} function of $\theta$, so the
sign of the angle label is immaterial and both methods produce the same
normalized curve. This evenness is also the source of the inherent
$\pm\theta$ ambiguity of magnitude-only correlation: the rotation can only be
recovered up to sign from the correlation alone. This ambiguity is already
resolved downstream in FS2D, since for every candidate rotation the 2D
translation correlation (Sec.\,\ref{sec:2DTranslation}) is evaluated and the
correct sign is selected by translation peak height or by an initial guess.

\paragraph{Two normalized-away differences.}
Two details differ between the raw (pre-normalization) outputs of the two
methods, both of which vanish after the $[0,1]$ normalization that precedes peak
detection:
\begin{itemize}
    \item \textbf{The $m=0$ DC term.} The full-SO(3) correlation includes the
          $m=0$ contribution, which is a real constant offset
          $P_{0}=\sum_{l}\hat{f}^{1}_{l0}\overline{\hat{f}^{2}_{l0}}$. The direct
          method explicitly skips $m=0$. Since $P_{0}$ is constant over $\theta$,
          it shifts the curve uniformly and disappears under min/max
          normalization.
    \item \textbf{The per-degree Wigner normalization.} The SO(3) coefficient
          combine in Eq.\,\ref{eq:so3coef} carries a per-degree factor
          $2\pi\sqrt{2/(2l+1)}$, which is compensated by the
          $\pi/(B\!\cdot\!N)$ normalization of the inverse SO(3) transform.
          The direct inner product in Eq.\,\ref{eq:Pm} does not include this
          factor. End to end, both methods produce the same $\mathbb{L}^{2}$
          correlation shape, so after normalization the curves match. Only if
          raw correlation values or per-degree weights are consumed downstream
          of the normalization does this factor need to be accounted for.
\end{itemize}

\subsubsection{What is no longer recomputed}
\label{sec:direct_dropped}

When the direct method is used, the following steps of the full-SO(3) path are
dropped entirely:
\begin{itemize}
    \item The SO(3) coefficient combine
          $\hat{f}^{1}_{lm},\hat{f}^{2}_{lm}\!\to\!\hat{C}^{l}_{mm'}$ with its
          Wigner-$d$ structure and per-degree normalization
          (Eq.\,\ref{eq:so3coef}).
    \item The inverse SO(3) Fourier transform, i.e.\ the $O(N^{3})$ step that
          produces the 3D correlation $C(\alpha,\beta,\gamma)$. This is the
          dominant cost of the full-SO(3) path.
    \item The SO(3)-specific Wigner-$d$ seminaive tables needed only for the
          inverse transform. The $\mathbb{S}^{2}$ forward tables used by
          \texttt{FST\_semi\_memo} are still required and unchanged.
    \item The $\beta$ Gauss--Legendre sampling and the extraction of the
          $\beta\approx 0$ slice, replaced by the exact analytic $\beta=0$
          reduction.
    \item The 2D-to-1D averaging over $(\alpha,\gamma)$ pairs with the same
          combined angle (Eq.\,\ref{eq:so3_collapse}), replaced by a 1D curve
          produced natively.
\end{itemize}

The following steps are \emph{unchanged} and shared by both methods: the 2D FFT
magnitude and phase, the FFT-shift and global-max normalization, the radial
$\mathbb{S}^{2}$ resampling, CLAHE and the Hamming window, both forward
spherical transforms \texttt{FST\_semi\_memo}, the rotation peak detection, and
the entire downstream pipeline (re-rotation of the scan, the 2D translation
correlation of Sec.\,\ref{sec:2DTranslation}, and the covariance estimation).
The translation half of FS2D is completely unaffected by the choice of rotation
method.

\paragraph{Complexity.}
Replacing the inverse SO(3) transform and the Wigner-$d$ combine by the direct
$P_{m}$ inner product and the 1D Fourier-series evaluation reduces the
asymptotic cost of the rotation step from
$O(N^{3})$ to $O(N^{2}\log N + N B)$ with $B=N/2$, while producing a normalized
1D rotation curve that is analytically equivalent to the full-SO(3) result and,
on the $\beta=0$ slice, strictly more exact.
